% ICA 2026 -- macros required by the corrected draft that
% scripts/ica26_build_paper_macros.py does not yet emit.
%
% This file is a TEMPORARY BRIDGE. Fold each definition into the generator and
% delete this file. Everything here that carries a value is DERIVED from numbers
% already in results_macros.tex by arithmetic, not measured afresh; the two
% \ResultPending entries need a short evaluation script (see below).
%
% \providecommand is used throughout, so once the generator emits a macro of the
% same name, the generator's definition (loaded first) wins and this file
% becomes a no-op.

% --------------------------------------------------------------------------- %
% 1. Calibration degradation ratio.
%
% Replaces the hand-typed claim "between 1.5 and 1.9 times", which was wrong at
% both ends: the actual ratios of the seed-mean cross-domain ECE are 1.48, 1.55
% and 1.94.
%
% Generator line (ratio of the seed-mean ECE, matching how ECE is reported):
%   ece_ratio[m] = mean_ece[m]["xd_T"] / mean_ece[m]["xd"]
%
% Preferred, stricter version: compute the ratio per seed and emit
%   \EceRatioMin / \EceRatioMax over all nine model-seed cases. If you do that,
%   change the sentence in Section 7 to quote the range instead of three values.
% --------------------------------------------------------------------------- %
\providecommand{\EceRatioRn}{1.48}
\providecommand{\EceRatioEb}{1.55}
\providecommand{\EceRatioMn}{1.94}

% --------------------------------------------------------------------------- %
% 2. Largest across-seed standard deviation of the fitted in-domain temperature.
%
% Replaces the hand-typed "never exceeds 0.008".
% Generator line:
%   temp_sd_max = max(sd_of_fitted_T[m] for m in backbones)   # PlantVillage
% --------------------------------------------------------------------------- %
\providecommand{\TempSdMax}{0.0076}

% --------------------------------------------------------------------------- %
% 3. Trivial-baseline reference for the cross-domain evaluation set.
%
% NEEDED: a reader cannot tell whether 0.26 accuracy over 21 classes is good,
% bad or trivial without one. This is the single cheapest addition in the paper
% and requires no training at all.
%
% Generator sketch (operates on the frozen cross-domain manifest only):
%   y = labels of the 1951 shared-class images
%   maj = most frequent label
%   acc  = (y == maj).mean()
%   f1   = sklearn.metrics.f1_score(y, np.full_like(y, maj),
%                                   average="macro", labels=present_classes)
%
% Report the majority-class predictor rather than uniform random: it is the
% stronger and therefore more honest floor under a skewed label distribution.
% --------------------------------------------------------------------------- %
\providecommand{\AccXdMajority}{\ResultPending}
\providecommand{\MacroFXdMajority}{\ResultPending}

% --------------------------------------------------------------------------- %
% 4. OPTIONAL but high value: in-domain performance restricted to the 21 shared
%    classes.
%
% This is the direct answer to Limitation 1 (label spaces differ). It needs no
% new training: take the SAVED PlantVillage test logits, apply the same
% canonical-class summation and renormalisation used cross-domain, and score
% only the PlantVillage test images whose class falls in the shared space.
% The resulting drop is then 21-way to 21-way, i.e. domain shift with task
% difficulty held fixed.
%
% If you generate these, add one paragraph to Section 6 and downgrade
% Limitation 1 from a threat to a quantified decomposition. Reviewers ask for
% exactly this.
% --------------------------------------------------------------------------- %
\providecommand{\MacroFPvInSharedRn}{\ResultPending}
\providecommand{\MacroFPvInSharedEb}{\ResultPending}
\providecommand{\MacroFPvInSharedMn}{\ResultPending}
\providecommand{\RelDropFSharedRn}{\ResultPending}
\providecommand{\RelDropFSharedEb}{\ResultPending}
\providecommand{\RelDropFSharedMn}{\ResultPending}
